\documentclass[12pt,a4paper]{article}
\usepackage[utf8]{inputenc}
\usepackage[T1]{fontenc}
\usepackage[italian]{babel}
\usepackage{lmodern}
\usepackage[a4paper,margin=2.5cm]{geometry}
\usepackage{setspace}
\usepackage{hyperref}
\usepackage{enumitem}
\usepackage{graphicx}

\setstretch{1.1}
\hypersetup{
    colorlinks=true,
    linkcolor=black,
    urlcolor=blue
}

\title{Classificazione binaria di segnali PCG con focus su preprocessing}
\author{
    \textbf{Artificial Intelligence from Engineering to Arts - aa 2025/2026}
    \\
    Gianluca Farinaccio - 548013 \\
    Alexis Osvaldo Hidalgo Martinez - 549484 \\
    \\
    \url{https://github.com/osvaldomar144/pcg-heart-sound-pipelines.git}
}

\begin{document}
\maketitle

\begin{abstract}
Questo progetto affronta la classificazione binaria di fonocardiogrammi (PCG) in classi
healthy/unhealthy, usando dati derivati dalla challenge PhysioNet/CinC 2016. L'obiettivo
principale è valutare in modo sistematico l'impatto del preprocessing prima dell'addestramento di
una CNN. Sono state confrontate pipeline modulari basate su trim, filtraggio band-pass e
normalizzazione RMS, seguite da trasformata Mel e compressione logaritmica. I risultati mostrano
che il preprocessing influisce in modo sostanziale sul compromesso tra accuracy globale e capacità
di riconoscere la classe minoritaria unhealthy. Il contributo principale del lavoro è quindi
metodologico: confronto riproducibile tra pipeline e analisi dei limiti di un baseline CNN su
dataset sbilanciato.
\end{abstract}

\section{Introduzione}
La classificazione automatica dei suoni cardiaci è un problema rilevante per screening precoce e
triage, soprattutto in contesti con risorse limitate. I segnali PCG sono però complessi da
modellare perché sono spesso affetti da rumore ambientale, artefatti di acquisizione e variabilità
fisiologica intra/inter-paziente. A questo si aggiunge lo sbilanciamento tra classi, che può
spostare il comportamento del modello verso la classe maggioritaria.

Nel progetto è stata costruita una pipeline end-to-end composta da preprocessing audio
configurabile, conversione in rappresentazione tempo-frequenza (Mel spectrogram) e classificazione
con CNN. La priorità è stata posta sul preprocessing, perché nella pratica i modelli PCG sono molto
sensibili a come il segnale viene filtrato, normalizzato e uniformato temporalmente.

\section{Stato dell'arte}
La challenge PhysioNet/CinC 2016 resta il benchmark storico più usato per classificazione
normal/abnormal su PCG. In generale, la letteratura riporta risultati sensibilmente migliori di
quelli ottenuti in questo progetto baseline, ma tali prestazioni derivano quasi sempre da pipeline
più complesse.

Le soluzioni di punta non si basano su una semplice CNN applicata a Mel-spectrogrammi: integrano
tipicamente preprocessing avanzato, segmentazione cardiaca esplicita delle fasi del battito e
strategie di riduzione del rumore più sofisticate. Inoltre, oltre alle CNN classiche, vengono
impiegati modelli sequenziali e ibridi, ad esempio LSTM/BiLSTM o combinazioni tra più blocchi
architetturali.

In sintesi, il confronto con lo stato dell'arte suggerisce che questo lavoro rappresenta un buon
punto di partenza metodologico, ma che per colmare il gap prestazionale sono necessari
preprocessing e modelli più avanzati.

\section{Metodi e tecnologie}
\subsection{Dataset}
Il dataset era inizialmente disponibile come unico blocco di training. Per costruire un protocollo
sperimentale completo, questo insieme è stato suddiviso in train, validation e test con
frazionamento stratificato 80\%-15\%-5\%, mantenendo una distribuzione healthy/unhealthy il più
coerente possibile tra i tre split.

I dati di base sono file audio \texttt{.wav} campionati a 2000 Hz, con durata variabile tra le
diverse registrazioni. Il dataset è inoltre eterogeneo per contesto di acquisizione: i punti di
auscultazione possono cambiare da file a file, ma questa informazione non è riportata in modo
strutturato nei metadati usati nel progetto. Anche la classe \texttt{unhealthy} è etichettata in
forma aggregata: segnala la presenza di anomalia, ma non specifica in modo sistematico la
patologia cardiaca sottostante.

Dal punto di vista statistico il problema resta sbilanciato. L'Imbalance Ratio, definito come
\(N(\text{healthy}) / N(\text{unhealthy})\), è circa 3.87: in pratica la classe healthy è quasi
quattro volte più rappresentata della classe unhealthy. Questo significa che un classificatore
banale che predice sempre healthy può ottenere un'accuracy apparentemente alta, senza essere però
utile in termini clinici.

In training, questo sbilanciamento tende a favorire la classe maggioritaria durante
l'ottimizzazione, con rischio di recall più basso sulla classe minoritaria e aumento dei falsi
negativi su unhealthy. Per questo la sola accuracy non basta: la valutazione deve includere anche
macro-F1, recall per classe e analisi della confusion matrix.

\subsection{Pipeline di preprocessing implementate}
Le pipeline sono definite in \texttt{configs/pipelines.json} con parametri globali condivisi
(\texttt{target\_sr=2000}, \texttt{chunk\_seconds=8.13}, \texttt{n\_fft=512},
\texttt{hop\_length=128}, \texttt{n\_mels=128}, banda Mel 20--300 Hz). Le varianti effettivamente
usate nel confronto sono tre: \texttt{trim}, \texttt{bandpass} e \texttt{bandpass\_norm\_rms}.

In tutte le pipeline il file audio viene caricato in mono, poi uniformato temporalmente con
\texttt{trim} alla durata target. Nella variante \texttt{bandpass} viene applicato un Butterworth
passa-banda (20--300 Hz, ordine 4, con \texttt{filtfilt}) per attenuare componenti fuori banda
cardiaca. Nella variante \texttt{bandpass\_norm\_rms} si aggiunge anche la normalizzazione RMS per
ridurre la variabilità di volume tra registrazioni. La rappresentazione finale è ottenuta con
\texttt{mel}, seguita da trasformazione \texttt{log1p} per comprimere la dinamica dello
spettrogramma e stabilizzare l'input alla rete.

I parametri delle pipeline non sono stati scelti in modo arbitrario, ma derivano da analisi
iterative dei segnali e dai vincoli imposti dall'input richiesto dalla CNN. In particolare, banda
del filtro, durata del chunk e parametri tempo-frequenza sono stati definiti cercando un
compromesso tra robustezza generale, preservazione dell'informazione utile e riduzione del rischio
di overfitting su questo specifico dataset.

Con i parametri scelti (\texttt{n\_mels=128}, \texttt{chunk\_seconds=8.13},
\texttt{hop\_length=128}), l'uscita di ogni pipeline è un log-mel spettrogramma di dimensione
128x128, salvato in scala di grigi.

\paragraph{Esempio qualitativo sano/non-sano}
Per visualizzare qualitativamente l'effetto del preprocessing, nelle Figure
\ref{fig:trim-log-pair} e \ref{fig:bnr-log-pair} è mostrata la stessa coppia di file
healthy/unhealthy dopo lo step \texttt{log}, rispettivamente per le pipeline \texttt{trim} e
\texttt{bandpass\_norm\_rms}. In entrambe le immagini il pannello sinistro è healthy e quello
destro è unhealthy.

\begin{figure}[htbp]
    \centering
    \includegraphics[width=\textwidth]{sanity/pipeline_trim_20260213_153137/04_log.png}
    \caption{Coppia healthy/unhealthy con pipeline \texttt{trim} (output log-mel)}
    \label{fig:trim-log-pair}
\end{figure}

\begin{figure}[htbp]
    \centering
    \includegraphics[width=\textwidth]{sanity/pipeline_bandpass_norm_rms_20260213_153230/06_log.png}
    \caption{Coppia healthy/unhealthy con pipeline \texttt{bandpass\_norm\_rms} (output log-mel)}
    \label{fig:bnr-log-pair}
\end{figure}

L'analisi visiva evidenzia, in entrambe le pipeline, una differenza coerente tra le due classi: il
campione healthy mostra una struttura temporale più ricca e una distribuzione energetica che si
estende maggiormente verso bande medio-basse, mentre il campione unhealthy presenta energia più
concentrata nelle basse frequenze e pattern più sparsi nel tempo. Nel confronto diretto tra
\texttt{trim} e \texttt{bandpass\_norm\_rms}, su questo esempio le differenze globali sono
moderate, ma la seconda pipeline produce mappe leggermente più pulite e con contrasto più
regolare, coerentemente con l'effetto combinato di filtraggio passa-banda e normalizzazione RMS.

\subsection{CNN - Architettura e parametri di training}
Il classificatore è una ResNet18, inizializzata con pesi pretrained ImageNet e adattata alla
classificazione binaria sostituendo il layer finale. Gli spettrogrammi vengono convertiti in
immagini grayscale replicate a 3 canali, ridimensionate a 224x224 e normalizzate.

Sono stati eseguiti tre addestramenti, uno per ogni pipeline sopra proposta. L'addestramento è
stato eseguito per 10 epoche, batch size 32, ottimizzatore Adam e loss CrossEntropy. Inoltre è
presente \texttt{WeightedRandomSampler} perché, come descritto sopra, il dataset presenta un forte
sbilanciamento, con Imbalance Ratio pari a 3.87.

\subsection{Stack tecnologico}
Per la realizzazione del progetto descritto si è utilizzato Python come linguaggio di scripting,
\texttt{numpy}/\texttt{scipy}/\texttt{librosa}/\texttt{soundfile} per l'elaborazione audio,
\texttt{matplotlib} per visualizzazione e reportistica, \texttt{torch} e \texttt{torchvision} per
training ed evaluation, e \texttt{scikit-learn} per le metriche di classificazione. L'intero
flusso è orchestrato da script CLI dedicati per preprocessing, training, sanity check e valutazione
su test set.

\section{Risultati}
La valutazione riporta i tre run \texttt{model\_best} delle pipeline \texttt{bandpass},
\texttt{bandpass\_norm\_rms} e \texttt{trim}. Tutti i confronti fanno riferimento allo stesso test
set (324 campioni: 258 healthy e 66 unhealthy).

\subsection*{Confronto delle metriche globali}
La Tabella~\ref{tab:metriche-globali} riassume il confronto principale su accuracy, precision,
recall e F1 in versione weighted.

\begin{table}[htbp]
    \centering
    \caption{Confronto metriche globali tra pipeline}
    \label{tab:metriche-globali}
    \begin{tabular}{|l|c|c|c|c|}
        \hline
        \textbf{Pipeline} & \textbf{Accuracy} & \textbf{Precision weighted} & \textbf{Recall weighted} & \textbf{F1 weighted} \\
        \hline
        \texttt{bandpass} & 0.7531 & 0.7418 & 0.7531 & 0.7470 \\
        \hline
        \texttt{bandpass\_norm\_rms} & 0.7222 & 0.8461 & 0.7222 & 0.7489 \\
        \hline
        \texttt{trim} & 0.5926 & 0.7682 & 0.5926 & 0.6318 \\
        \hline
    \end{tabular}
\end{table}

La metrica F1 weighted è particolarmente utile in questo contesto perché tiene conto della
numerosità reale delle classi nel test set, evitando che il risultato globale sia distorto da classi
con pochi campioni. Allo stesso tempo, resta più informativa della sola accuracy perché combina
precision e recall in un unico indicatore.

La pipeline \texttt{bandpass} ottiene la migliore accuracy, mentre \texttt{bandpass\_norm\_rms}
mostra il miglior F1 weighted. La pipeline \texttt{trim} resta la meno efficace sul piano globale.

\subsection*{Focus sulla classe unhealthy}
Poiché la classe unhealthy è minoritaria e clinicamente rilevante, la
Tabella~\ref{tab:metriche-unhealthy} mostra precision, recall e F1 specifiche su questa classe.

\begin{table}[htbp]
    \centering
    \caption{Metriche sulla classe unhealthy}
    \label{tab:metriche-unhealthy}
    \begin{tabular}{|l|c|c|c|}
        \hline
        \textbf{Pipeline} & \textbf{Precision unhealthy} & \textbf{Recall unhealthy} & \textbf{F1 unhealthy} \\
        \hline
        \texttt{bandpass} & 0.3793 & 0.3333 & 0.3548 \\
        \hline
        \texttt{bandpass\_norm\_rms} & 0.4143 & 0.8788 & 0.5631 \\
        \hline
        \texttt{trim} & 0.2963 & 0.7273 & 0.4211 \\
        \hline
    \end{tabular}
\end{table}

Il dato più rilevante è il recall unhealthy: \texttt{bandpass\_norm\_rms} riduce in modo netto i
falsi negativi rispetto alle altre pipeline, mentre \texttt{bandpass} privilegia la classe healthy
e perde molti campioni unhealthy.

\subsection*{Supporto visuale}
Le Figure~\ref{fig:cm-normalizzate} e \ref{fig:global-metrics-plots} mostrano rispettivamente le
confusion matrix normalizzate e i grafici delle metriche globali per le tre pipeline.

\begin{figure}[htbp]
    \centering
    \includegraphics[width=0.32\textwidth]{metrics/pipeline_banpass_model_best_test/confusion_matrix_normalized.png}
    \hfill
    \includegraphics[width=0.32\textwidth]{metrics/pipeline_bnr_model_best_test/confusion_matrix_normalized.png}
    \hfill
    \includegraphics[width=0.32\textwidth]{metrics/pipeline_trim_model_best_test/confusion_matrix_normalized.png}
    \caption{Confusion matrix normalizzate: \texttt{bandpass} (sinistra), \texttt{bandpass\_norm\_rms}
    (centro), \texttt{trim} (destra)}
    \label{fig:cm-normalizzate}
\end{figure}

\begin{figure}[htbp]
    \centering
    \includegraphics[width=0.32\textwidth]{metrics/pipeline_banpass_model_best_test/global_metrics.png}
    \hfill
    \includegraphics[width=0.32\textwidth]{metrics/pipeline_bnr_model_best_test/global_metrics.png}
    \hfill
    \includegraphics[width=0.32\textwidth]{metrics/pipeline_trim_model_best_test/global_metrics.png}
    \caption{Grafici metriche globali: \texttt{bandpass} (sinistra), \texttt{bandpass\_norm\_rms}
    (centro), \texttt{trim} (destra)}
    \label{fig:global-metrics-plots}
\end{figure}

In sintesi, i risultati mostrano due comportamenti distinti: \texttt{bandpass} è la scelta migliore
se l'obiettivo è massimizzare accuracy complessiva, mentre \texttt{bandpass\_norm\_rms} è la scelta
più solida quando si vuole privilegiare la rilevazione della classe unhealthy (recall più alto e
F1 weighted migliore).

\section{Conclusioni}
Il progetto ha prodotto una base sperimentale solida e riproducibile per confrontare pipeline di
preprocessing in classificazione PCG binaria. Il risultato principale è che il preprocessing
modifica in modo sostanziale il trade-off tra specificità sulla classe healthy e sensibilità sulla
classe unhealthy: non esiste quindi una pipeline ``migliore'' in assoluto, ma una scelta legata
alla metrica obiettivo e al contesto di utilizzo.

La principale evidenza metodologica è che, in presenza di sbilanciamento, accuracy e val\_acc
possono essere ottimistiche se non accompagnate da metriche per-classe, F1 weighted e analisi
degli errori. Un limite importante del lavoro è anche la quantità di dati disponibile: il numero di
campioni, in particolare per la classe unhealthy, riduce la capacità di generalizzazione del
modello e aumenta la sensibilità alle specificità del dataset usato.

Come sviluppi futuri, risultano quindi prioritari: un confronto multi-seed/cross-validation,
strategie di data augmentation mirate per aumentare variabilità e bilanciare meglio le classi, e la
validazione dell'intera pipeline su dataset esterni più ampi ed eterogenei. In parallelo, resta
rilevante valutare modelli ibridi con componente sequenziale o attenzione e calibrare la soglia
decisionale in funzione del costo clinico dei falsi negativi.

\section*{Riferimenti web (fonti primarie)}
\begin{enumerate}[leftmargin=*,label=\arabic*.]
    \item PhysioNet Challenge 2016 (overview ufficiale): \url{https://physionet.org/content/challenge-2016/}
    \item Zabihi et al., Heart Sound Classification from Unsegmented Phonocardiograms Using Deep Learning (CinC 2016): \url{https://www.cinc.org/archives/2016/pdf/236-175.pdf}
    \item Springer et al., Logistic Regression-HSMM-Based Heart Sound Segmentation (IEEE TBME 2016): \url{https://pubmed.ncbi.nlm.nih.gov/26340769/}
\end{enumerate}

\end{document}
